\subsection{2026-09-22 --- Budget augmentation: +0.020 from the structural insight}
\label{sec:budget}

Section~\ref{sec:oracle} established that the entire gap from 0.83 to 0.98 lies in one number per
unit-week, and that a naive reformulation replacing the pointwise model performed worse than it.
This entry records three further attempts and the one that worked.

\subsubsection*{Unit of analysis.}
The count problem lives at \textbf{(case, unit)}: 40{,}908 rows, one per unit-week. This
granularity resolves a limitation carried since Section~\ref{sec:feature-themes}: in the
(case, hour) layout a static per-unit attribute is a constant column with zero information, so
$p_{\max}$, $p_{\min}$, \texttt{startcost\_const}, \texttt{penstock} and merit rank were unusable.
With \emph{unit} as a row dimension rather than a column dimension they become ordinary features
that vary across rows. The class mix at this level is 27\% all-off, 41\% mixed, 32\% all-on.

\subsubsection*{Three attempts.}

\begin{center}
\begin{tabular}{llr}
\toprule
Attempt & Form & micro-AUC (2022) \\
\midrule
1 & flat count regression, hard allocation      & 0.7801 \\
2 & 3-class $+$ count regression, hard allocation & 0.8257 \\
  & \emph{pointwise model, for reference}       & \emph{0.8347} \\
3 & \textbf{pointwise $+$ budget offset (logit scale)} & \textbf{0.8548} \\
\bottomrule
\end{tabular}
\end{center}

\finding{\textbf{Decomposing the count target helped, but replacing the pointwise model did not.}
Splitting into all-off / mixed / all-on and regressing the count only on mixed unit-weeks cut count
MAE from \textbf{43.6h to 34.6h} and lifted the allocation score from 0.7801 to 0.8257. Both
remain below the pointwise model, because hard allocation turns a count error into a step function
that ranks badly, whereas the pointwise model degrades smoothly.}

\finding{\textbf{Used as an \emph{augmentation} rather than a replacement, the same information is
worth $+0.020$.} Adding $\lambda \cdot (\widehat{\text{count}} - \text{spread rank})/168$ to the
pointwise model's logits gives 0.8548 against a 0.8347 baseline. The pointwise model ranks hours
well within a week but has no representation of \emph{how many} hours a unit should run; the
budget margin supplies exactly that, while the logit form preserves the existing ranking instead
of overriding it.}

\finding{\textbf{The blend weight transfers across years, unlike the stage-2 gate.} $\lambda$
selected on 2021 (=4) yields $+0.0200$ on 2022, against $+0.0240$ for the 2022-optimal
$\lambda{=}8$ --- only 0.004 lost to honest selection. Both years show the same broad shape (rise,
flat peak in the 4--8 range, decline), so the choice is insensitive. This contrasts with the
per-unit stacking gate in Section~\ref{sec:stacked}, which was coin-flip across years and therefore
unusable.}

\subsubsection*{A methodological error, caught by a sanity check.}
The first blend attempt rank-normalised each unit's predictions independently before combining.
The sweep included $\lambda{=}0$, which must by construction reproduce the pointwise baseline; it
returned 0.7872 against 0.8347, exposing the error. \textbf{Micro-averaged AUC pools every
(row, unit) cell into a single ranking, so per-column rank transforms destroy the cross-unit
calibration} the metric depends on --- the information that Vesle\_Kjela is almost always on while
Vinje\_G3 is usually off. Blending must happen on the logit scale. Recorded because the invalid
version looked entirely reasonable and produced plausible numbers; only the fixed point in the
sweep revealed it.

\subsubsection*{Remaining headroom.}
Count MAE is 34.6h against a target mean of 87h, and the oracle ceiling is 0.9765. The augmented
model captures $+0.020$ of the $+0.142$ available, so the count model is still the binding
constraint. The 2022 confusion matrix also shows a distribution-shift problem: all-on unit-weeks
are 32\% of training but only 9\% of 2022 (the price-crisis year, Section~\ref{sec:eda-round-2}
finding 9), and only 176 of 474 are recovered. A count model robust to that shift --- or trained
with recency weighting --- is the obvious next target.
